\section{Schopenhauer on the Will}

Schopenhauer recognized that space, time, and causality are required for any experience and understanding of the natural world, or the world of phenomena. Knowledge, then, consists in understanding the sufficient reason for things, that is, the causal relationships between them. He also recognized, following Kant, that the phenomenal world is not self-explanatory or self-sufficient: there must be something beyond the natural, phenomenal world. Thus, he claimed, there is the noumenal, which he named the ``Will''.

However, as a rationalist and philosopher, this Will is beyond knowledge, in the sense defined above. As unknowable, it can only be described as blind, purposeless, and unpredictable. Otherwise, we would know its sufficient reason, and it would be just another thing among other things in the natural world, and no longer the one ultimate reality and cause.

From a psychological point of view, the Will manifests as our drives, instincts, desires, and so on. Since it is purposeless, so, too, are our drives, instincts, and desires. This is the cause of pain and suffering, so his solution lies in the rejection of the Will. The negative consequence is that it leads to a life of quietism and resignation.

From a metaphysical point of view, we would say he recognizes the faculty of reason, but rejects the higher form of knowledge called ``intellect'', which knows transphenomenal reality directly and intuitively.


\flrightit{Posted on 2009-08-11 by Cologero}

